% Project Christopher -- 보고서(report/main.tex)와 입문 가이드(report/guide/main.tex)의 공용 전처리부.
% 엔진: XeLaTeX + kotex. 글꼴은 TeX Live에 들어 있는 것만 쓴다(은글꼴, Latin Modern).

\usepackage[a4paper,margin=25mm,headheight=14pt]{geometry}
% \usepackage{fontspec}
\usepackage[hangul]{kotex}
% \setmainhangulfont{UnBatang.ttf}[BoldFont=UnBatangBold.ttf,ItalicFont=UnBatang.ttf,ItalicFeatures={FakeSlant=0.17},BoldItalicFont=UnBatangBold.ttf,BoldItalicFeatures={FakeSlant=0.17}]
% \setsanshangulfont{UnDotum.ttf}[BoldFont=UnDotumBold.ttf,ItalicFont=UnDotum.ttf,BoldItalicFont=UnDotumBold.ttf]
% \setmonohangulfont{UnDotum.ttf}[BoldFont=UnDotumBold.ttf,ItalicFont=UnDotum.ttf,BoldItalicFont=UnDotumBold.ttf]

\usepackage{amsmath,amssymb}
\usepackage{graphicx}
\usepackage{pdfpages}     % 회로도 부록: KiCad PDF를 쪽 단위로 싣기
\usepackage{xcolor}
\usepackage{booktabs,array,tabularx,multirow}
\usepackage{caption,subcaption}
\usepackage{enumitem}
\usepackage{setspace}
\setstretch{1.25}
\setlength{\parskip}{0.45em}
\setlength{\parindent}{0pt}
\setlist{itemsep=0.15em,topsep=0.3em}
\captionsetup{font=small,labelfont=bf,width=0.92\linewidth}

% ---------- 색 ----------
\definecolor{cblue}{HTML}{1D4ED8}
\definecolor{cred}{HTML}{B91C1C}
\definecolor{cgreen}{HTML}{15803D}
\definecolor{corange}{HTML}{C2410C}
\definecolor{cpurple}{HTML}{6D28D9}
\definecolor{cgray}{HTML}{6B7280}
\definecolor{clight}{HTML}{F3F4F6}
\definecolor{cbluebg}{HTML}{DBEAFE}
\definecolor{credbg}{HTML}{FEE2E2}
\definecolor{cgreenbg}{HTML}{DCFCE7}
\definecolor{corangebg}{HTML}{FFEDD5}
\definecolor{cyellowbg}{HTML}{FEF9C3}

% ---------- 그림 ----------
\usepackage{tikz}
\usepackage{pgfplots}
\usepackage{pgfplotstable}
\pgfplotsset{compat=1.18}
\usepackage{tikz-timing}
\usepackage{circuitikz}   % 회로도(릴레이 접점, 스위치)
\usetikzlibrary{arrows.meta,positioning,calc,fit,backgrounds,shapes.geometric,shapes.misc,decorations.pathreplacing,matrix,patterns}
\tikzset{
  >={Stealth[length=2.2mm]},
  every picture/.style={font=\small},
  blk/.style={draw,rounded corners=1.5pt,minimum height=8mm,minimum width=16mm,align=center,fill=white,inner sep=3pt},
  blkb/.style={blk,fill=cbluebg,draw=cblue},
  blkg/.style={blk,fill=cgreenbg,draw=cgreen},
  blko/.style={blk,fill=corangebg,draw=corange},
  blkr/.style={blk,fill=credbg,draw=cred},
  blkn/.style={blk,fill=clight,draw=cgray},
  reg/.style={draw,minimum height=10mm,minimum width=9mm,fill=cyellowbg,align=center,inner sep=2pt},
  cloud/.style={draw,rounded corners=4mm,minimum height=9mm,minimum width=20mm,fill=cbluebg,draw=cblue,align=center,inner sep=3pt},
  ltr/.style={circle,draw,minimum size=6.5mm,inner sep=0pt,fill=white,font=\small\ttfamily},
  note/.style={font=\footnotesize,align=center,text=cgray},
  wire/.style={line width=0.5pt},
  live/.style={line width=1.1pt,draw=cblue},
  rcoil/.style={draw,fill=cyellowbg,minimum width=9mm,minimum height=6mm,inner sep=1pt,font=\scriptsize,align=center},
  shaft/.style={line width=2.2pt,draw=cgray},
}

% ---------- 상자 ----------
\usepackage[most]{tcolorbox}
% 아직 쓸 수 없는 부분(자리 표시)
\newtcolorbox{todobox}[1]{enhanced,breakable,colback=clight,colframe=cgray,boxrule=0.6pt,
  borderline={0.6pt}{0pt}{cgray,dashed},frame hidden,arc=1.5pt,left=3mm,right=3mm,top=2mm,bottom=2mm,
  fonttitle=\bfseries\sffamily\small,coltitle=black,colbacktitle=clight,
  title={\textcolor{cred}{[작성 예정]}\ #1}}
% 이 프로젝트에서 실제로 겪은 일
\newtcolorbox{lesson}[1]{enhanced,breakable,colback=corangebg!45,colframe=corange,boxrule=0.6pt,arc=1.5pt,
  left=3mm,right=3mm,top=2mm,bottom=2mm,fonttitle=\bfseries\sffamily\small,title={실제 사례: #1}}
% 핵심 요약
\newtcolorbox{keybox}[1][]{enhanced,breakable,colback=cbluebg!45,colframe=cblue,boxrule=0.6pt,arc=1.5pt,
  left=3mm,right=3mm,top=2mm,bottom=2mm,fonttitle=\bfseries\sffamily\small,#1}

% ---------- 코드 ----------
\usepackage{listings}
\lstdefinelanguage{Rust}{
  morekeywords={as,break,const,continue,crate,else,enum,extern,false,fn,for,if,impl,in,let,loop,match,mod,move,mut,
    pub,ref,return,self,Self,static,struct,super,trait,true,type,unsafe,use,where,while,u8,u16,u32,u64,usize,f64,bool,Option,Some,None,Vec},
  sensitive=true,morecomment=[l]{//},morecomment=[s]{/*}{*/},morestring=[b]"}
\lstdefinelanguage{SV}[]{Verilog}{
  morekeywords={logic,always_comb,always_ff,typedef,enum,int,genvar,localparam,parameter,package,endpackage,import,
    automatic,function,endfunction,unique,priority,bit,return},
  morecomment=[l]{//}}
\lstdefinestyle{code}{basicstyle=\ttfamily\footnotesize,keywordstyle=\color{cblue}\bfseries,
  commentstyle=\color{cgreen},stringstyle=\color{corange},columns=fullflexible,keepspaces=true,
  frame=single,rulecolor=\color{cgray!50},backgroundcolor=\color{clight!60},xleftmargin=2mm,xrightmargin=2mm,
  framesep=2mm,aboveskip=0.8em,belowskip=0.6em,showstringspaces=false,breaklines=true,tabsize=2,
  numbers=left,numberstyle=\tiny\color{cgray},numbersep=6pt,captionpos=b}
\lstset{style=code}
\renewcommand{\lstlistingname}{코드}
\newcommand{\code}[1]{\texttt{#1}}

% ---------- 참고문헌 ----------
\usepackage[backend=biber,style=numeric,sorting=none,maxbibnames=4]{biblatex}

% ---------- 기타 ----------
\usepackage[colorlinks=true,linkcolor=cblue!70!black,citecolor=cgreen!60!black,urlcolor=cblue!70!black]{hyperref}
\newcommand{\term}[2]{\textbf{#1}(#2)}       % 용어(원어)
\newcommand{\Scr}{\mathrm{Scr}}
\newcommand{\Stk}{S}
